\documentclass[11pt]{article}

\usepackage[margin=1in]{geometry}
\usepackage{amsmath,amssymb}
\usepackage{natbib}
\usepackage{graphicx}
\usepackage{booktabs}
\usepackage{xcolor}
\usepackage{hyperref}
\usepackage{cleveref}

\newcommand{\GMSL}{\text{GMSL}}
\newcommand{\GMST}{\text{GMST}}

\title{Budget-Consistent Sea Level Component Projections\\
  with Observationally Constrained Ice Sheet Uncertainty}

\author{B.\ Minchew}
\date{Working Draft --- \today}

\begin{document}
\maketitle

% ====================================================================
% TARGET: GMD (alongside FACTS) or Earth's Future
% LENGTH: ~6000-7000 words
% FIGURES: 6-8 main
% TIMELINE: Submittable 2-3 months after Paper A
% ====================================================================

\begin{abstract}
% ~200 words
Probabilistic sea level rise projections typically sum independent draws from component-level distributions (thermal expansion, glaciers, ice sheets, terrestrial water storage), discarding the constraint that components must sum to the observed total.
We demonstrate a post-processing method that enforces this budget constraint on existing FACTS v1.0 component distributions, producing revised projections that are jointly consistent with satellite-era observations.
Resolving Antarctica into West Antarctic (WAIS), East Antarctic (EAIS), and Antarctic Peninsula (APIS) sub-regions, we show that the budget constraint acts as an asymmetric filter: it preferentially eliminates process-model runs that show WAIS mass gain under warming---runs inconsistent with the observed sea level budget---truncating the lower tail of the WAIS distribution and shifting the median upward.
For the critical WAIS component, we separate the projection into a budget-constrained committed trajectory and a deep-uncertainty marine ice-sheet instability (MISI) component, parameterized by an activation probability $\pi(t)$.
The committed component is anchored by observations; the MISI component extends the upper tail using structured expert judgment and high-end process-model output.
The resulting projection distributions are provided for all SSP scenarios under multiple $\pi$ assumptions, offering budget-consistent projections with transparent treatment of the WAIS deep uncertainty that dominates high-end sea level risk.
\end{abstract}


% ====================================================================
\section{Introduction}
\label{sec:intro}
% ~600 words
% ====================================================================

% Key points:
% 1. FACTS (Kopp et al. 2023) produces component-level SLR projections by
%    independently sampling each component's distribution. The total is the
%    sum of independent draws. The budget identity (components sum to total)
%    is not enforced during sampling.
%
% 2. This independence assumption wastes information. During the satellite era,
%    total GMSL is observed with high precision. The budget constraint creates
%    a collider mechanism (Pearl 2009) that induces posterior dependence between
%    components: well-observed components (thermal expansion) tighten the
%    residual that constrains poorly observed components (WAIS).
%
% 3. For WAIS specifically, the budget constraint addresses a known problem:
%    ISMIP6 models (Seroussi et al. 2020) largely fail to produce marine
%    ice-sheet instability. Most models show Antarctic mass gain under warming.
%    These runs are inconsistent with the observed budget and should receive
%    lower weight.
%
% 4. But the budget constraint alone is insufficient for WAIS projections.
%    It constrains the historical trajectory but cannot constrain a future
%    threshold process (MISI) that has not yet fully activated. We therefore
%    separate the WAIS projection into a committed component (budget-constrained)
%    and a MISI component (physically motivated, parameterized).
%
% 5. This paper provides a method, not just a diagnosis. The output is a set
%    of revised, budget-consistent projection distributions that users can
%    adopt directly.


% ====================================================================
\section{Method}
\label{sec:method}
% ~2000 words
% ====================================================================

\subsection{Component decomposition}
\label{sec:components}
%
% We decompose GMSL into 7 components:
% TE, GL, GrIS, WAIS, EAIS, APIS, LWS
%
% Table 1: Component definitions, forcing variables, observational data sources.
%
% All FACTS AIS modules (emulandice, larmip, bamber19, deconto21) provide
% WAIS/EAIS/APIS via full_sample_components. This sub-continental resolution
% is necessary because WAIS, EAIS, and APIS have distinct physics:
% - WAIS: marine, CDW-forced, MISI-susceptible
% - EAIS: land-based, SMB-dominated, near steady state
% - APIS: atmospheric warming, hydrofracture

\subsection{Budget constraint via importance sampling}
\label{sec:budget}
%
% Given N independent draws from the 7-component FACTS prior:
%   {r_i^(n)}_{i=1}^7, n = 1,...,N
%
% Compute importance weights from the budget likelihood:
%   w^(n) = exp(-0.5 sum_t [(sum_i r_i^(n)(t) - r_obs(t))^2 / sigma^2(t)])
%
% where r_obs(t) is the observed total GMSL rate and sigma^2(t) includes
% observation uncertainty + budget closure tolerance.
%
% The weighted ensemble approximates the posterior:
%   p({r_i} | r_obs) propto p({r_i}) * L(r_obs | {r_i})
%
% Multiple observational constraints can be applied:
% - Global budget: sum of all 7 = total GMSL (Frederikse, altimetry)
% - Antarctic sub-budget: WAIS + EAIS + APIS = total AIS (IMBIE, Frederikse, Horwath)
% - Component-level: individual IMBIE, Argo, GlaMBIE constraints where available
%
% Present equation for combined log-likelihood with all constraint levels.

\subsection{The budget constraint as a physical filter on ISMIP6}
\label{sec:filter}
%
% Key result, developed in detail:
%
% The ISMIP6 ensemble (Seroussi et al. 2020, 2024) contains models that show
% Antarctic mass gain under warming. In the FACTS emulandice/AIS module, these
% correspond to samples with rate_WAIS < 0 (mass gain) or rate_WAIS ~ 0.
%
% When the budget constraint is applied:
% 1. Compute the implied total for each sample.
% 2. Samples where WAIS contributes near-zero or negative, combined with
%    the well-constrained positive contributions from TE, GL, GrIS, produce
%    a total that is too low relative to observed GMSL.
% 3. These samples receive low importance weights.
% 4. The WAIS posterior has its lower tail truncated. The median shifts up.
%
% Quantify: what fraction of emulandice/AIS WAIS samples are effectively
% eliminated (weight < 0.01 * max_weight)? What is the pre- vs post-constraint
% WAIS median and 17-83% range?
%
% This is an observationally grounded reweighting of the ISMIP6 ensemble.
% It does not require knowing why the models fail (though the n=3 rheology
% bias and ocean forcing deficiencies are known mechanisms). It only requires
% that the models be consistent with the observed total.
%
% Figure 1: WAIS prior vs posterior distribution at 2050 and 2100.
% Show the lower tail truncation.

\subsection{WAIS committed + MISI mixture model}
\label{sec:mixture}
%
% The budget-constrained posterior addresses the low tail but does not
% generate the upper tail. The ISMIP6 models cannot produce MISI. The
% budget constraint cannot constrain a process that has not activated.
%
% We separate the WAIS projection into two components:
%
% p(r_WAIS(t)) = (1 - pi(t)) * p_committed(r_WAIS(t)) 
%                + pi(t) * p_MISI(r_WAIS(t))
%
% --- p_committed ---
% The budget-constrained WAIS trajectory, extrapolated forward using the
% posterior rate and acceleration at the end of the observational period.
% This is the "no MISI" scenario: WAIS continues to accelerate at
% approximately the observed rate, driven by ongoing CDW intrusion and
% grounding-line retreat on prograde/flat bed slopes.
%
% Specification:
% - Initial rate r_0 and acceleration dr/dt at 2020 from the budget-constrained
%   posterior (not from IMBIE alone — the budget shifts the estimate).
% - Forward model: rate(t) = r_0 + (dr/dt)(t - t_0) + 0.5*(d2r/dt2)(t-t_0)^2
%   with the second derivative set to zero (conservative) or estimated from
%   the budget-constrained posterior curvature.
% - Uncertainty: propagated from the posterior on r_0 and dr/dt.
%
% --- pi(t) ---
% The probability that MISI has activated by time t. This is the deep-
% uncertainty parameter. It is NOT constrained by the budget.
%
% Specification options (present all, recommend transparent parameterization):
%
% Option 1: Transparent parameterization.
%   Present projections for pi(2100) in {0, 0.05, 0.10, 0.25, 0.50}.
%   pi(t) ramps linearly from 0 at present to pi(2100) at 2100.
%   The reader supplies their own belief. This is the primary presentation.
%
% Option 2: Process-model calibration.
%   Seroussi et al. (2024): 30-40% of extended ISMIP6 models show widespread
%   WAIS collapse by 2300 under high emissions. Infer pi(t) from the fraction
%   of models that have collapsed by each time step. This gives a process-model-
%   informed pi that can be compared to the expert-judgment-informed pi.
%
% Option 3: Expert judgment.
%   Extract pi implicitly from Bamber et al. (2019): compare the Bamber WAIS
%   distribution to the committed trajectory. The excess probability in the
%   Bamber upper tail relative to p_committed is attributable to the experts'
%   implicit belief about MISI. Invert to get pi.
%
% Option 4: Observational early-warning.
%   Use the Rosier et al. (2021) tipping-point analysis for PIG. PIG may
%   already be past its instability threshold. Thwaites is approaching it.
%   Assign pi_PIG ~ 0.7-1.0 (likely committed) and pi_Thwaites as a function
%   of CDW forcing trajectory. This physically grounds pi in specific glacier
%   systems. [More appropriate for Paper C than Paper B — flag for later.]
%
% --- p_MISI ---
% The WAIS rate distribution conditional on MISI having activated.
%
% Specification:
% - Source 1: Upper tail of Bamber et al. (2019) WAIS distribution.
%   Take the Bamber samples exceeding the committed 83rd percentile.
%   These represent the experts' assessment of what happens with MISI/MICI.
%
% - Source 2: DeConto et al. (2021) / deconto21 WAIS output.
%   Under SSP5-8.5, deconto21 produces substantially more WAIS loss than
%   emulandice. This excess is attributable to the MICI mechanism in the
%   DeConto model. Extract the deconto21-minus-committed residual.
%
% - Source 3: Getraer & Morlighem (2025) n=4 scaling.
%   ISMIP6 with n=4 produces 32% more ASE mass loss by 2100, 70% by 2300.
%   Apply this scaling factor to the committed trajectory to get a
%   "rheology-corrected MISI" estimate.
%
% Use Source 1 (Bamber) as primary. Sources 2 and 3 as sensitivity checks.
% Fit a shifted log-normal to the MISI-conditional rate distribution.
%
% Figure 2: Schematic of the mixture model.
% Show p_committed (narrow, budget-anchored) and p_MISI (wide, heavy-tailed)
% and the mixture at different pi values.


% ====================================================================
\section{Results}
\label{sec:results}
% ~1500 words
% ====================================================================

\subsection{Budget constraint tightens component distributions}
%
% Table 2: Variance ratio rho_j for all 7 components at 2050 and 2100
% under SSP2-4.5, using emulandice modules, combined constraint.
%
% Key results:
% - WAIS: largest variance reduction (rho ~ 0.6-0.75)
% - TE: minimal reduction (rho ~ 0.95) — already well-constrained
% - EAIS, APIS: moderate reduction
%
% Figure 3: 7x7 posterior correlation matrix at 2100.
% Off-diagonal entries show the collider-induced dependence.
% WAIS-TE correlation should be strongly negative.
% WAIS-EAIS correlation negative (Antarctic sub-budget collider).

\subsection{Asymmetric filtering of ISMIP6 WAIS distribution}
%
% Figure 4: Prior vs posterior WAIS distribution.
% Two panels: emulandice/AIS and larmip/AIS.
% Show the lower tail truncation and median shift.
% Annotate: fraction of prior samples effectively eliminated.
%
% Key number: X% of emulandice WAIS samples are inconsistent with the
% observed budget. The budget-constrained median is Y mm/yr higher than
% the unconstrained median.

\subsection{WAIS projections with MISI uncertainty}
%
% Figure 5: WAIS projection distributions at 2100.
% Multiple panels or overlaid densities:
% - pi = 0 (committed only, budget-constrained)
% - pi = 0.10
% - pi = 0.25
% - pi = 0.50
% Under SSP2-4.5 and SSP5-8.5.
%
% Show how the mixture extends the upper tail without affecting the
% budget-constrained center/lower tail.
%
% Table 3: WAIS projection quantiles at 2100 for each SSP x pi combination.
% Format: median (17-83) [5-95].

\subsection{Total GMSL projections}
%
% Figure 6: Total GMSL projection fan chart.
% Combine the budget-constrained non-WAIS components with the WAIS mixture
% to produce total GMSL projections.
% Show for pi = 0.10 (moderate) and pi = 0.25 (substantial).
% Compare to FACTS unconstrained total.
%
% Table 4: Total GMSL at 2050, 2100, 2150 for each SSP × pi.
% IPCC-comparable format (meters, relative to 1995-2014).
%
% Key result: the budget-constrained projection with moderate MISI probability
% (pi = 0.10-0.25) produces a total that:
% - Is similar to FACTS in the median (budget correction offsets MISI addition)
% - Has a narrower 17-83% range at 2050 (budget constraint dominates)
% - Has a wider 83-95% range at 2100 (MISI tail extends beyond FACTS)
% - Has better-calibrated component decomposition (no compensating errors)


% ====================================================================
\section{Discussion}
\label{sec:discussion}
% ~800 words
% ====================================================================

\subsection{What the budget constrains and what it cannot}
%
% Clear statement of epistemics:
% - The budget anchors WAIS during the historical period. It eliminates
%   model trajectories inconsistent with observations.
% - It cannot constrain future threshold crossings. MISI is parameterized,
%   not predicted.
% - The separation into committed + MISI is not arbitrary: it tracks the
%   distinction between what the observational record can inform and what
%   requires additional physical/expert judgment.

\subsection{Relationship to existing frameworks}
%
% - FACTS: this method is a post-processing step applicable to FACTS output.
%   It does not replace FACTS; it refines it. The code and revised distributions
%   will be made available in a format compatible with FACTS workflows.
%
% - Bamber et al. (2019): the expert judgment distributions implicitly contain
%   MISI. Our method makes the MISI component explicit and separates it from
%   the budget-constrained committed component. This transparency is an
%   advantage: the user sees exactly what comes from observations and what
%   comes from expert judgment.
%
% - Paper A (DOLS projections): Paper A provides the thermodynamic projection;
%   Paper B provides the component decomposition with budget consistency and
%   MISI treatment. They are complementary. Paper A's DOLS residual provides
%   an independent check on the ice-sheet contribution diagnosed by Paper B.

\subsection{The n = 3 rheology bias}
%
% Brief discussion of how the Martin et al. (2025) n=3 initialization bias
% compounds the WAIS underestimate. The budget constraint partially corrects
% this (it eliminates runs that are too low regardless of cause), but it
% cannot fully correct a structural model deficiency. The n=4 scaling from
% Getraer & Morlighem (2025) provides a physics-based correction that could
% be applied to the committed trajectory. We present this as a sensitivity
% analysis, not a primary result.

\subsection{Implications for AR7}
%
% The revised distributions are directly usable. They provide:
% 1. Budget-consistent component projections (no compensating errors).
% 2. WAIS projections with transparent MISI treatment.
% 3. Projection tables in IPCC-comparable format.
% 4. Code to apply the budget constraint to future FACTS releases.


% ====================================================================
\section{Data and Code Availability}
% ====================================================================
% Revised projection distributions: Zenodo deposit (CSV, NetCDF)
% Budget constraint code: GitHub repository
% FACTS post-processing pipeline: documented, open-source
% All input data: publicly available


% ====================================================================
% SUPPLEMENTARY MATERIALS
% ====================================================================
% SM1: Importance sampling methodology details
% SM2: Budget constraint specification (sigma_budget, observational products)
% SM3: Antarctic sub-budget constraint (Frederikse, Horwath cross-validation)
% SM4: WAIS committed trajectory specification (rate, acceleration extraction)
% SM5: MISI distribution fitting (Bamber tail extraction, log-normal fit)
% SM6: pi(t) specification options (transparent, process-model, expert)
% SM7: Sensitivity to AIS module choice
% SM8: Sensitivity to sigma_budget
% SM9: ESS diagnostics
% SM10: n=4 rheology scaling sensitivity
%
% Fig S1: AIS cross-validation (IMBIE vs Frederikse vs Horwath)
% Fig S2: Variance ratio heatmap (7 components × time × SSP)
% Fig S3: Constraint level comparison (global only vs combined vs AIS-only)
% Fig S4: Within-Antarctica information flow
% Fig S5: WAIS prior vs posterior for all 4 AIS modules
% Fig S6: pi(t) sensitivity for total GMSL
% Fig S7: n=4 rheology scaling effect on WAIS projections
% Fig S8: ESS as function of sigma_budget
% Fig S9: Budget residual time series (1900-2020)
%
% Table S1: Full variance ratios (all components × SSPs × constraint types)
% Table S2: Full 7x7 correlation matrices (prior and posterior)
% Table S3: WAIS filtering statistics (fraction eliminated, median shift)
% Table S4: Full projection tables (decadal, all SSP × pi combinations)
% Table S5: Sensitivity to sigma_budget


\bibliographystyle{agsm}
\bibliography{slr_references}

\end{document}
